\chapter{Reproducibility}
\label{app:repro}

Everything claimed in this report can be regenerated from the
repository.\footnote{\url{https://github.com/witold-andelie/quant-foundation-gat}}
This appendix records the environment, the commands, and the
mapping from experiment to archived artifact.

\section*{Environment}

Python 3.13 with \texttt{torch 2.12.0} and \texttt{torch\_geometric}
(installed via the project's \texttt{[gnn]} extra); \texttt{PYTHONPATH}
set to \texttt{src}. All headline experiments ran on CPU --- a
deliberate choice after the device study of \cref{sec:eq-variance} ---
with the device recorded per experiment. The equity track fetches
market data through \texttt{yfinance}; the energy track runs on the
synthetic generator by default and on live ENTSO-E data when an API
token is supplied via the \texttt{ENTSOE\_API\_KEY} environment
variable. The warehouse and marts require \texttt{duckdb},
\texttt{dbt-core}, and \texttt{dbt-duckdb}; the GAT pipeline itself does
not.

\section*{Tests}

The full suite is 115 tests; 55 of them cover the capstone modules
(graph construction, training primitives, propagators, factor
providers, both track runners, attention analysis, and the leakage
controls). The leakage controls --- shuffle-label negative control and
planted-signal positive control --- are pinned to CPU for determinism
and run on every change. \texttt{py -m pytest} from the repository root
executes everything.

\section*{Experiments and artifacts}

Each experiment entry in the log (\texttt{docs/gat\_experiment\_log.md},
entries E1--E14) names its run script and its archived outputs. The run
scripts are one-off drivers kept in an untracked scratch area outside
the published repository; each log entry records the configuration they
ran --- model variant, seeds, hyperparameters, device --- and the CLI
commands below drive the same pipelines. The outputs --- per-run
diagnostics, matrix, seed, hyperparameter-grid, winner-validation,
energy, and attention CSV/JSON files, plus five attention figures ---
are date-prefixed under \texttt{docs/results/}. The canonical trained
equity weights are tracked at
\path{data/warehouse/gat_equity.pt}. Every quantitative table in
this report is generated mechanically from those archived CSVs; no
result number in the text was typed by hand from memory.

\section*{Command reference}

The equity pipeline:
\begin{verbatim}
quant-alpha gat-equity --graph {static,dynamic}
                       --retrain {single,walk-forward}
                       --loss {ic,mse}
                       --device {cpu,cuda,auto}
                       --persist
\end{verbatim}
runs end to end and, with \texttt{--persist}, writes the four warehouse
tables of \cref{sec:pl-loop}. The energy alpha runner is
\texttt{run\_gat\_energy} (synthetic or \texttt{source="entsoe"}); the
forecasting ladder is \texttt{quant-alpha energy-forecast --source
\{synthetic,entsoe\} --k 24} with a \texttt{--raw-path} Parquet cache
for repeatable real-data runs. The marts rebuild with
\texttt{dbt build --profiles-dir .} inside \texttt{dbt\_quant\_alpha/}
(21 models and tests) and \texttt{dbt\_energy\_alpha/} (12). The
report's build script regenerates the tables and figures from the
repository's artifacts and reference data and then runs
\texttt{pdflatex} and \texttt{biber}.

\section*{What reproducibility means here}

As established in \cref{sec:eq-variance}, bit-exact replication across
devices is not attainable for CUDA-adjacent training stacks and is not
claimed. Rerunning any experiment with the recorded seeds on a CPU of
the same class reproduces the reported distributions --- means and
standard deviations over five or more seeds --- and every individual
archived run is identified by seed, device, configuration, and date in
its artifact file.
